\begin{frame}
  \onslide<1->{
	\frametitle{The Static Estimation Problem} 
	\framesubtitle{Derivation of the Kalman Gain}
}
\begin{itemize}
  \onslide<2->{
 \item Given: A known measurement vector $z=(z_1,z_2,\ldots
 ,z_m)$ that is believed to be causally related to the
 state vector $x=(x_1,\ldots ,x_n)$.
}
  \onslide<3->{
 \item Find: The "best" \textit{\textbf{estimate}} of the state as a function of the
 measurements: $\hat{x}=g(z)=g(z_1,\ldots ,z_m)$.
}
\end{itemize}
\end{frame}


\begin{frame}
  \onslide<1->{
	\frametitle{The Static Estimation Problem} 
	\framesubtitle{Bayesian Derivation of Kalman Measurement Update}
}
\begin{itemize}
  \onslide<2->{
 \item Ignore state dynamics for now -- hence the term "static"
}
  \onslide<3->{
 \item Disregard stochastic differential equations and don't have to use mean square calculus
}
  \onslide<4->{
 \item Concerned only with statistical problem of estimating a vector random variable $x$ given a finite set of
 measurements
}
\end{itemize}
\end{frame}

\begin{frame}
  \onslide<1->{
	\frametitle{Static Estimation Problem} 
	\framesubtitle{Why should we care?}
}
\begin{itemize}
  \onslide<2->{
 \item Solution to this problem will introduce concepts central to traditional Kalman filtering
}
  \onslide<3->{
 \item Solution here nearly identical to that given by Kalman himself in his 1960 paper; result is
 now classical
}
  \onslide<4->{
 \item We'll derive the \textbf{Kalman gain} matrix $K$ that is used to weight the measurement residual in
 updating the a priori state estimate and also used to update the a priori estimation error covariance matrix $P$
}
\end{itemize}
\end{frame}
\begin{frame}
  \onslide<1->{
	\frametitle{The Static Estimation Problem} 
	\framesubtitle{Definition of Optimality}
}
\begin{itemize}
  \onslide<2->{
 \item What do we mean by "best" estimate?
}
  \onslide<3->{
 \item Need to define a measure of "goodness", i.e., a performance
 measure and maximize it.
}
  \onslide<4->{
 \item Actually, it turns out to be more convenient to define a
 measure of "badness" that we can minimize.
 \begin{itemize}
    \item This is called a "loss function" or "penalty function".
 \end{itemize}
}
\end{itemize}
\end{frame}
\begin{frame}
  \onslide<1->{
	\frametitle{The Static Estimation Problem} 
	\framesubtitle{Mean Squared Error}
}
\begin{itemize}
  \onslide<2->{
 \item Estimation Error: The difference between the truth and the
 estimate: $\tilde{x}=\hat{x}-x$
}
  \onslide<3->{
 \item Estimation error is bad
}
  \onslide<4->{
    \item Our estimators should try to minimize this error.
}
  \onslide<5->{
 \begin{itemize}
    \item At the very least, we should try to keep track of how
    well we're doing (i.e., should have an error covariance
    matrix).
 \end{itemize}
}
\end{itemize}
\end{frame}

\begin{frame}
  \onslide<1->{
	\frametitle{The Static Estimation Problem} 
	\framesubtitle{Admissible Loss Functions}
}
\begin{itemize}
  \onslide<2->{
    \item Nonnegative, real-valued function of the estimation
    error;
}
  \onslide<3->{
    \item Non-decreasing, i.e., the loss gets bigger as the
    estimation error gets bigger;
}
  \onslide<4->{
    \item Should be zero when there is no error, i.e., when we
    are certain;
}
  \onslide<5->{
    \item Should not be dependent on the sign of the error, i.e.,
    should be a symmetric function.
}
  \onslide<6->{
 \item Note: As a function of a random variable, the loss function
 has a random variable as its value.  Hence, we'll want to minimize the
 mean of the loss function.
}
\end{itemize}
\end{frame}


\begin{frame}
  \onslide<1->{
	\frametitle{The Static Estimation Problem} 
	\framesubtitle{Means Squared Error}
}
\begin{itemize}
  \onslide<2->{
 \item Estimation problem becomes: Find estimate $\hat{x}$ of $x$
 that minimizes the mean of the loss function $L$:\\
 $E\left[L(\tilde{x})\right]=E\left[L(\hat{x}-x)\right]$
}
  \onslide<3->{
 \item A common loss function is a quadratic function: the square of the norm of the error
 vector:  $L(\tilde{x})=\tilde{x}^\top\tilde{x}$ \\
 Often referred to as the "square error."
}
  \onslide<4->{
 \item Therefore, we want to minimize the mean:
 $E\left[L(\tilde{x})\right]=E\left[\tilde{x}^\top\tilde{x}\right]$\\
 i.e., we want to minimize the "mean square error."
}
  \onslide<5->{
 \item Note that $L$ is really a function of two variables:
 $L(\tilde{x})=L(\hat{x},x)$.
}
\end{itemize}
\end{frame}

\begin{frame}
  \onslide<1->{
	\frametitle{The Static Estimation Problem} 
	\framesubtitle{Problem Formulation}
}
\begin{itemize}
  \onslide<2->{
 \item Estimation problem becomes: Find estimate $\hat{x}$ of $x$
 that minimizes the integral:\\
\begin{equation}\label{lossfunc}
 E\left[\tilde{x}^\top\tilde{x}\right]=\underbrace{\int_{-\infty}^{+\infty}\cdots \int_{-\infty}^{+\infty}}_{n+m}\tilde{x}^\top\tilde{x}f(x,z)dx_1\cdots dx_n dz_1\cdots dz_m
\end{equation}

 where $f(x,z)$ is the joint probability density function for $x$
 and $z$.
}
  \onslide<3->{
 \item $x$ is assumed to be an $n$-vector, while $z$ is assumed to
 be an $m$-vector.
}
\end{itemize}
\end{frame}

\begin{frame}
  \onslide<1->{
	\frametitle{The Static Estimation Problem} 
	\framesubtitle{Solving the Problem}
}
\begin{itemize}
  \onslide<2->{
 \item But we can factor the joint p.d.f, through the conditional p.d.f:
    \[f(x,z)=f(x|z)f(z)\]

    where

    $f(z)=\underbrace{\int_{-\infty}^{+\infty}\cdots\int_{-\infty}^{+\infty}}_{n}f(x,z)dx_1\cdots dx_n$
    is the marginal p.d.f of $z$.
}
\end{itemize}
\end{frame}

\begin{frame}
  \onslide<1->{
	\frametitle{The Static Estimation Problem} 
	\framesubtitle{Solving the Problem}
}
\begin{footnotesize}
  \onslide<2->{
Therefore, we can modify
Eq.~(\ref{lossfunc}), our integral expression for
$E\left[\tilde{x}^\top\tilde{x}\right]$, as follows:
}
\begin{align*}
  \onslide<3->{
  E\left[\tilde{x}^\top\tilde{x}\right] & = \underbrace{\int_{-\infty}^{+\infty}\cdots
\int_{-\infty}^{+\infty}}_{n+m}\tilde{x}^\top\tilde{x}f(x,z)dx_1\cdots
dx_n dz_1\cdots dz_m\\
}
  \onslide<4->{
    & = \underbrace{\int_{-\infty}^{+\infty}\cdots
\int_{-\infty}^{+\infty}}_{n+m}\tilde{x}^\top\tilde{x}f(x|z)f(z)dx_1\cdots
dx_n dz_1\cdots dz_m\\
}
  \onslide<5->{
    & = \underbrace{\int_{-\infty}^{+\infty}\cdots
\int_{-\infty}^{+\infty}}_{m}f(z)\left[\underbrace{\int_{-\infty}^{+\infty}\cdots
\int_{-\infty}^{+\infty}}_{n}\tilde{x}^\top\tilde{x}f(x|z)dx_1\cdots
dx_n\right]dz_1\cdots dz_m
}
\end{align*}
  \onslide<6->{
Note that we have pulled $f(z)$ out of the $x$-integral, because
$f(z)$ does not depend on $x$.
}
\end{footnotesize}
\end{frame}

\begin{frame}
  \onslide<1->{
	\frametitle{The Static Estimation Problem} 
	\framesubtitle{Solving the Problem}
}
\begin{itemize}
  \onslide<2->{
 \item Notice that we now
only need to minimize the integral within the brackets, i.e.,

$E\left[\tilde{x}^\top\tilde{x}\right]$ is minimized if and only if\\

$\underbrace{\int_{-\infty}^{+\infty}\cdots
\int_{-\infty}^{+\infty}}_{n}\tilde{x}^\top\tilde{x}f(x|z)dx_1\cdots
dx_n$ is minimized.
}
  \onslide<3->{
Remember that we are trying to find the value of $\hat{x}$ that
will cause this loss function to be minimized.
}
\end{itemize}
\end{frame}

\begin{frame}
  \onslide<1->{
	\frametitle{The Static Estimation Problem} 
	\framesubtitle{Solving the Problem}
}
\begin{itemize}
  \onslide<2->{
\item Let's expand $\tilde{x}^\top\tilde{x}:$
}
  \onslide<3->{
\begin{align*}
  \qloss & = \normsq{(\hat{x}-x)} \\
    & =\hat{x}^\top\hat{x}-\hat{x}^\top x - x^\top \hat{x} + x^\top
   x\\
   & =\normsq{\hat{x}}+\dprod{x}{x}- 2\dprod{\hat{x}}{x}
\end{align*}
}
\end{itemize}
\end{frame}

\begin{frame}
  \onslide<1->{
	\frametitle{The Static Estimation Problem} 
	\framesubtitle{Solving the Problem}
}
\begin{itemize}
  \onslide<2->{
\item Thus, we expand the integral to be minimized accordingly:\\
}
  \onslide<3->{
\begin{multline}
  \minfint{n}\normsq{\tilde{x}}f(x|z)dx_1\cdots dx_n  = \normsq{\hat{x}}\\
    +\minfint{n}\normsq{x}f(x|z)dx_1\cdots dx_n\\
    - 2\hat{x}\minfint{n}xf(x|z)dx_1\cdots dx_n 
\end{multline}
}
\end{itemize}
\end{frame}
\begin{frame}
  \onslide<1->{
	\frametitle{The Static Estimation Problem} 
	\framesubtitle{Solving the Problem}
}
\begin{small}
\begin{itemize}
  \onslide<2->{
\item  This last expansion is justified because $\normsq{\hat{x}}$ is a
function of the $z$ alone:  $\hat{x}=g(z_1,\ldots ,z_m)$.  The
equation of interest has the form:
 
}
  \onslide<3->{
\begin{equation}\label{loss_integrand}
\minfint{n}\normsq{\tilde{x}}f(x|z)dx_1\cdots dx_n =
f_1(\hat{x})+f_2(z) -2f_3(\hat{x},z)
\end{equation}
}
  \onslide<4->{
where $f_1(\hat{x})=\normsq{\hat{x}}$,
$f_2(z)=\minfint{n}\normsq{x}f(x|z)dx_1\cdots dx_n$ and
}
  \onslide<5->{
$f_3(\hat{x},z)= \hat{x}^\top\minfint{n} f(x|z)dx_1\cdots dx_n$.
}
\end{itemize}
\end{small}
\end{frame}

\begin{frame}
  \onslide<1->{
	\frametitle{The Static Estimation Problem} 
	\framesubtitle{Solving the Problem}
}
\begin{itemize}
  \onslide<2->{
\item   Now we need to minimize the right side of the last equation. \\ To
do so, we must take the gradient with respect to $\hat{x}$,\\ set it
equal to zero, and solve for $\hat{x}$ -- a typical min/max
problem.
}
  \onslide<3->{
\item   Recall that the gradient of the dot product is:\\
\begin{center}
$\pderiv{\dprod{y}{x}}{x}=y$ and also $\pderiv{\normsq{x}}{x}=2x$.
\end{center}
}
\end{itemize}
\end{frame}

\begin{frame}
  \onslide<1->{
	\frametitle{The Static Estimation Problem} 
	\framesubtitle{Solving the Problem}
}
\begin{small}
\begin{itemize}
  \onslide<2->{
\item Applying this to our problem, we note that the middle term of \\
Eq. ~(\ref{loss_integrand}), 
$f_2(z)$ is a function of $z$ \\
alone and thus vanishes when the gradient \\
 with respect to $\hat{x}$ is applied to it.
}
  \onslide<3->{
\item  We are thus left with the first and third terms only:\\

$\pderiv{\normsq{\hat{x}}}{\hat{x}}=2\hat{x}$ ~~~and \\

\begin{multline}
\pderiv{(\hat{x}^\top\minfint{n} x f(x|z)dx_1\cdots dx_n)}{\hat{x}} = \\
 \minfint{n} x f(x|z)dx_1\cdots dx_n
\end{multline}

}
\end{itemize}
\end{small}
\end{frame}

\begin{frame}
  \onslide<1->{
	\frametitle{The Static Estimation Problem } 
	\framesubtitle{Solving the Problem}
}
\begin{itemize}
  \onslide<2->{
\item  Thus, taking the gradient with respect to $\hat{x}$ of
Eq.~(\ref{loss_integrand}) we obtain \\

$0 = 2\hat{x}-2\minfint{n}xf(x|z)dx_1\cdots dx_n$ ~~~~ or\\

}
  \onslide<3->{
\begin{equation}\label{opt_est}
  \hat{x}=\minfint{n}xf(x|z)dx_1\cdots dx_n
\end{equation}
  
}
\end{itemize}
\end{frame}

\begin{frame}
  \onslide<1->{
	\frametitle{The Static Estimation Problem} 
	\framesubtitle{The Solution}
}
\begin{itemize}
  \onslide<2->{
\item This solution is just the conditional expectation for $x$,
conditioned on the $z$:\\  
}
  \onslide<3->{
  $\hat{x}=\minfint{n}xf(x|z)dx_1\cdots dx_n=E\left[ x|z\right]$.

}
\end{itemize}
\end{frame}

\begin{frame}
  \onslide<1->{
	\frametitle{The Static Estimation Problem} 
	\framesubtitle{Assuming Gaussian Stats}
}
\begin{itemize}
  \onslide<2->{
\item   Now let's see what this looks like with Gaussian statistics for
$f(x|z):$\\

}
  \onslide<3->{
\begin{equation}\label{cond_mean}
  E\left[ x|z\right]=m_x+P_{xz}P^{-1}_{zz}(z-m_z),
\end{equation}
where $m_x, m_z$ are the means of $x$ and $z$, respectively, and
$P_{xz}$ is the cross-covariance matrix between $x$ and $z$, while
$P_{zz}$ is the covariance matrix for $z$ alone.
}
 \onslide<4->{
  \item Not that this is the matrix analogue of the scalar case for two Gaussians:
  $\expect{X|Z}=\expect{X}+\dfrac{\rho_{xz}\sigma_x}{\sigma_z}(z-\expect{Z})$
 }
\end{itemize}
\end{frame}
\begin{frame}
  \onslide<1->{
	\frametitle{The Static Estimation Problem} 
	\framesubtitle{Assuming Gaussian Stats}
}
\begin{itemize}
  \onslide<2->{
\item   In other words, we have for Eq.~(\ref{cond_mean}):
}
  \onslide<3->{
\begin{align*}
E\left[ x|z\right]&=m_x+P_{xz}P^{-1}_{zz}(z-m_z) \text{, where}\\
m_x &= E[ x ]\\
m_z &= E[ z ]\\
P_{xz}&=E[(x-m_x)(x-m_z)^\top]\\
P_{zz}&=E[(z-m_z)(z-m_z)^\top]
\end{align*}
}
\end{itemize}
\end{frame}

\begin{frame}
  \onslide<1->{
	\frametitle{The Static Estimation Problem} 
	\framesubtitle{The Solution}
}
\begin{itemize}
  \onslide<2->{
\item  Eq.~(\ref{cond_mean}) conveys the essence of the Kalman update algorithm.  
}
  \onslide<3->{
\item  It will give us a new estimate of
the state vector x conditioned on the actually observed
measurements and all associated statistics.
}
  \onslide<4->{
\item  Now consider again the relationship between x and z that we have
previously assumed:  $z=Hx+v$
}
  \onslide<5->{
Remember that $x$ and $z$ are statistically independent of each
other.
}
\end{itemize}
\end{frame}

\begin{frame}
  \onslide<1->{
	\frametitle{The Static Estimation Problem} 
	\framesubtitle{The Solution}
}
\begin{small}
\begin{itemize}
  \onslide<2->{
\item   Since $m_z=E[z]=E[Hx+v]=H\cdot E[x]+E[v]=Hm_x+m_v$, we have
}
  \onslide<3->{
\begin{align*}
P_{xz} &= E[(x-m_x)(z-m_z)^\top]\\
    &= E[(x-m_x)\{Hx+v - (Hm_x+m_v) \}^\top]\\
    &= E[(x-m_x)\{(Hx-Hm_x) + (v-m_v) \}^\top]\\
    &= E[(x-m_x)\{H(x-m_x) + (v-m_v) \}^\top]\\
    &= E[(x-m_x)\{(x-m_x)^\top H^\top + (v-m_v)^\top \}]\\
    &= E[(x-m_x)(x-m_x)^\top H^\top] + E[(x-m_x)(v-m_v)^\top]\\
    &= E[(x-m_x)(x-m_x)^\top H^\top]\\
\end{align*}
}
\end{itemize}
because $x$,$z$ independent $\Rightarrow E[(x-m_x)(v-m_v)^\top]=0$.
\end{small}
\end{frame}

\begin{frame}
  \onslide<1->{
	\frametitle{The Static Estimation Problem} 
	\framesubtitle{Translating into "Kalman Notation"}
}
\begin{itemize}
  \onslide<2->{
\item   But, $E[(x-m_x)(x-m_x)^\top] = P_{xx}$ by definition. \\

}
  \onslide<3->{
\item    Thus
\begin{align*}
P_{xz} &= E[(x-m_x)(x-m_x)^\top H^\top]\\
        &= P_{xx}H^\top
\end{align*}

}
  \onslide<4->{
\item   A similar derivation shows that
\begin{align*}
P_{zz} &= HP_{xx}H^\top+P_{vv}\\
\end{align*}

}
  \onslide<5->{
\item  Later we'll see this in the form of $HP_{xx}H^\top+R$ since
$P_{vv}=R$ by definition.
}
\end{itemize}
\end{frame}

\begin{frame}
  \onslide<1->{
	\frametitle{The Static Estimation Problem} 
	\framesubtitle{Translating into "Kalman Notation"}
}
\begin{itemize}
  \onslide<2->{
\item  Inserting these last two results into Eq.~(\ref{cond_mean}) yields
 
}
  \onslide<3->{
\begin{align*}
  E\left[ x|z\right]&=m_x+P_{xz}P^{-1}_{zz}(z-m_z)\\
        &= m_x+P_{xx}H^\top[HP_{xx}H^\top+P_{vv}]^{-1}(z-m_z)\\
        &= m_x+P_{xx}H^\top[HP_{xx}H^\top+P_{vv}]^{-1}(z-Hm_x)
\end{align*}
}
  \onslide<4->{
\item  This last step results because \\
$m_z=E[z]=E[Hx+v]=H\cdot E[x]+E[v]=Hm_x+0=Hm_x$\\
by definition.
}
\end{itemize}
\end{frame}


\begin{frame}
  \onslide<1->{
	\frametitle{The Static Estimation Problem} 
	\framesubtitle{Translating into "Kalman Notation"}
}
\begin{itemize}
  \onslide<2->{
\item The Kalman Gain:   Eq.~(\ref{cond_mean}) now becomes the central result
of classical Kalman Filter theory:
}
  \onslide<3->{
\begin{equation}\label{update_eq}
  \hat{x}=E[x|z]=m_x + K(z-Hm_x)
\end{equation}
where $K=P_{xx}H^\top[HPH^\top + P_{vv}]^{-1}$ is the \textbf{Kalman gain}.
}
\end{itemize}
\end{frame}
